% Modello canonico per le sezioni di esercizi dei capitoli teorici MQF.
% Compatibile con Scientific WorkPlace 5.5.
% Sostituire NN con il numero del capitolo e usare identificativi univoci.
% L'ambiente exercise e' definito in MQF_Manuale_Master.tex.

% ============================================================
% Esercizi svolti
% ============================================================

\section{Esercizi svolti}
\label{sec:capNN-esercizi-svolti}

\begin{exercise}[Titolo sintetico]
\label{ex:capNN-identificativo-svolto}
Presentare il contesto, i dati, le unita' di misura e le ipotesi.

\begin{enumerate}
\item Formulare o calcolare la prima quantita' richiesta.
\item Svolgere il controllo teorico o numerico pertinente.
\item Interpretare il risultato in termini economico-finanziari.
\end{enumerate}
\end{exercise}

\noindent
\textbf{Soluzione.}
\begin{enumerate}
\item Svolgimento motivato della prima richiesta.
\item Verifica coerente con le ipotesi dell'esercizio.
\item Interpretazione proporzionata ai risultati ottenuti.
\end{enumerate}

% Ripetere il blocco exercise + Soluzione per gli altri esercizi svolti.

% ============================================================
% Esercizi proposti
% ============================================================

\section{Esercizi proposti}
\label{sec:capNN-esercizi-proposti}

\begin{exercise}[Titolo sintetico]
\label{ex:capNN-identificativo-proposto}
Presentare il contesto, i dati, le unita' di misura e le ipotesi.

\begin{enumerate}
\item Formulare o calcolare la prima quantita' richiesta.
\item Svolgere il controllo teorico o numerico pertinente.
\item Interpretare il risultato in termini economico-finanziari.
\end{enumerate}

\smallskip
\textit{Risultati:} riportare soltanto i risultati essenziali per
l'autovalutazione, senza inserire lo svolgimento completo.
\end{exercise}

% Ripetere il blocco exercise per gli altri esercizi proposti,
% disponendoli secondo difficolta' progressiva.
